\section{Method}
\label{sec:method}

% CLAIM: C02
The method implements the transformations in
\cref{sec:problem-formulation} as explicit interfaces.  The estimator describes
the reference state supported by samples, the predictor assigns a future
reference time, the component selector defines the P/PV/PVA information
condition, the governor constructs an adjacent executable target, and the
follower emits the profile actually committed.  Plant feedback supplies the
next replanning-state candidate but is not silently substituted for the
reference posterior.

\subsection{Estimator and future-reference interfaces}

% CLAIM: C02,C05
An estimator consumes a position measurement carrying source and availability
times and returns \(\PosteriorState{\ell}{k}\).  Its represented time may lag
the control tick.  Missing, nonfinite, duplicate, or regressing samples are
handled by an explicit raise, hold, component-hold, or reset policy; a held
posterior remains stale.  For asynchronous axes, scalar posteriors are
propagated independently to the common control time under an explicitly
labelled constant-acceleration synchronization model.

% CLAIM: C02
The future-reference generator consumes a posterior but does not modify
estimator state.  Its output records the posterior source time, physical
prediction horizon, target time, availability, and whether the method is
causal.  Constant-velocity, constant-acceleration, constant-jerk, and local
polynomial models are instances of this interface.  Analytic future lookup is
kept as a separate offline, noncausal oracle; it is not presented as an online
predictor.

% CLAIM: C02
Equation~\eqref{eq:target-component-selection} then exposes P, PV, or PVA to
the downstream method.  This operation preserves position and timing while
zeroing omitted derivatives.  It changes the target-state information
condition, not the governor, follower, or plant identity.  Consequently, a
same-follower comparison must hold those later layers fixed.

\subsection{Authoritative current state}

Before planning, a single replanning state is selected and supplied to both
the governor and follower.  In \emph{previous-command} mode it is the last
committed endpoint; in \emph{measured} mode it is the available plant
measurement; and in \emph{hybrid} mode measurement feedback is selected when
command--measurement divergence exceeds the configured threshold.  The
selection, correction, and reason are recorded.  A measurement correction can
move the selected state outside the normal viability set, in which case the
normal invariant-preserving path is not asserted.

\subsection{One-step bounded-jerk executable-target governor}

% CLAIM: C09
For each axis, the governor selects one constant jerk over
\([\ControlTime{k},\ControlTime{k+1}]\) and constructs the endpoint using
\cref{eq:direct-a,eq:direct-v,eq:direct-p}.  With
\(x=[p,v,a]^{\mathsf T}\), this is
\begin{equation}
  x^+(j)=Ax+Bj,\qquad
  A=
  \begin{bmatrix}
    1&\DT&\tfrac12\DT^2\\
    0&1&\DT\\
    0&0&1
  \end{bmatrix},
  \quad
  B=
  \begin{bmatrix}
    \DT^3/6\\ \DT^2/2\\ \DT
  \end{bmatrix}.
  \label{eq:state-affine-in-jerk}
\end{equation}
The multi-axis construction applies this synchronized update independently
per axis under componentwise bounds.

Let \(\rho=\rho_{k,H}^{m}\) be the selected future-reference objective,
\(b=Ax\), and \(g=B\).  The implemented tracking preference is
\begin{equation}
  \mathcal L(j)=
  \sum_{c\in\{p,v,a\}}
    w_c\left(\frac{b_c+g_cj-\rho_c}{s_c}\right)^2
  +\lambda_j\left(\frac{j}{\JerkLimitSymbol}\right)^2
  +\lambda_{\Delta j}
    \left(\frac{j-j_{k-1}}{\JerkLimitSymbol}\right)^2,
  \label{eq:governor-objective}
\end{equation}
where
\(s=[\max(\VelocityLimitSymbol\DT,\epsilon),
\VelocityLimitSymbol,\AccelerationLimitSymbol]^{\mathsf T}\).
The first term tracks the requested state, while the remaining terms regularize
jerk magnitude and its change.  If \(H\neq1\), the first term deliberately
compares a next-cycle candidate with a farther-horizon objective; its residual
is therefore a look-ahead shaping objective, not same-time estimation error.

% CLAIM: C09
The segment-and-terminal jerk interval union is computed analytically.  It
intersects the jerk bound, endpoint acceleration and velocity bounds, any
interior velocity extremum, and terminal stopping viability.  A separate
postcheck verifies that a subsequent viable action exists.  The unconstrained
scalar quadratic preference is clipped to the analytic interval union.  If
numerical boundary effects invalidate that candidate, deterministic boundary
and zero-terminal-acceleration candidates are verified explicitly.  This
construction is a one-step, implementation-level selection rule; it is not a
globally optimal trajectory solution.
\Cref{app:governor-derivation} gives the interval construction and guarantee
boundary.

\begin{figure}[t]
  \centering
  \includegraphics[width=\linewidth]{governor_reachability.pdf}
  \caption{Illustrative endpoint slice for the one-step construction at
  \(v_k=a_k=0\), \(\Delta t=\SI{\ControlPeriodMS}{\milli\second}\), and the
  registered V/A/J limits.  The green curve shows only the constant-jerk
  endpoints inside the V/A box; the black point is the zero-jerk endpoint.
  The full acceptance rule additionally applies continuous segment checks,
  stopping viability, and next-step existence as defined in the text.}
  \label{fig:governor}
\end{figure}

% CLAIM: C09
Three outcomes are distinguished.  For a finite request and viable current
state, the normal path returns a verified constant-jerk endpoint.  For a
nonfinite request from a viable current state, the governor selects a verified
bounded-jerk recovery action and records an applied fallback.  If the selected
current state is already outside the viability conditions, the governor
integrates a finite, jerk-bounded best-effort emergency action; this branch is
labelled as emergency and does not inherit the normal guarantee.  No branch
implements an impossible state hold while acceleration or velocity dynamics
continue.

\subsection{Direct execution}

% CLAIM: C09
The direct follower is intended for an already executable target.  From the
selected current and proposed endpoint it computes
\begin{equation}
  j_{\mathrm{end}}=
  \frac{a_{k+1}^{\mathrm{target}}-a_k^{\mathrm{cur}}}{\DT}
  \label{eq:direct-endpoint-jerk}
\end{equation}
and independently reconstructs position and velocity using
\cref{eq:direct-v,eq:direct-p}.  Acceptance requires endpoint reconstruction,
continuous segment V/A/J legality, terminal stopping viability, and next-step
existence.  The implementation can additionally require
\(\FreeDuration(\CurrentState{k},\ExecutableTargetState{k+1})\leq\DT\) from a
separate ordinary-\Ruckig{} solve.  This named duration is an additional
diagnostic; it is not reused for a replacement endpoint.

When these checks pass, the executed command profile contains one segment with
duration \(\DT\), and
\(\CommandState{k+1}=\ExecutableTargetState{k+1}\).  A failed check invokes a
separately recorded one-step fallback constructed from the actual current
state.  The returned status distinguishes a requested fallback from one
actually committed and names the controller that generated the replacement.

\subsection{Ordinary \Ruckig{} prefix execution}

% CLAIM: C11
The \OrdinaryRuckig{} follower solves a synchronized state-to-state problem
from \(\CurrentState{k}\) to its requested endpoint.  It first obtains
\(\FreeDuration(\CurrentState{k},x^{\mathrm{req}})\) without a minimum-duration
constraint for request diagnostics.  A separate solve with the configured
minimum duration supplies the trajectory whose prefix over \([0,\DT]\) is
executed.  The prediction horizon never changes this minimum duration, and
neither quantity is a computation deadline.  A requested endpoint can have
\(\FreeDuration(\CurrentState{k},x^{\mathrm{req}})>\DT\) while the native
prefix endpoint remains the valid next-cycle command.

% CLAIM: C11
A native prefix can contain several constant-jerk segments.  Let
\(0=s_0<s_1<\cdots<s_M=\DT\), let
\(\delta_r=s_{r+1}-s_r\), and let \(j_r\) be the jerk on segment \(r\).
Starting from \(x_0=\CurrentState{k}\), the authoritative endpoint is
\begin{equation}
  x_{r+1}=\Phi_{\delta_r}(x_r,j_r),
  \qquad r=0,\ldots,M-1,
  \qquad
  \CommandState{k+1}=x_M.
  \label{eq:piecewise-prefix}
\end{equation}
Acceleration extrema occur at segment endpoints.  Velocity is also evaluated
at any interior \(\tau=-a_r/j_r\) that lies in
\((0,\delta_r)\), and jerk is checked on every segment.  Across multiple axes,
the union of exposed switch times defines a vector-valued piecewise profile.

% CLAIM: C11
The acceleration-difference quantity
\begin{equation}
  j_k^{\Delta a}=
  \frac{a_{k+1}^{\mathrm{cmd}}-a_k^{\mathrm{cur}}}{\DT}
  \label{eq:acceleration-difference-jerk}
\end{equation}
matches only the net acceleration change.  For a multi-segment prefix it is
not an internal jerk and generally fails to reconstruct terminal position and
velocity.  The compatibility field called command jerk therefore carries the
first actual profile jerk when available; the complete profile remains the
executable command.

When the binding exposes all brake, main, and acceleration segments, the
prefix is reconstructed and audited analytically, and its independently
integrated endpoint must match the native endpoint.  When exact segments are
not exposed, the implementation stores an explicitly inexact sampled state
grid augmented with accessible switch boundaries.  This supports endpoint and
sampled V/A auditing, but internal jerk remains unavailable rather than being
inferred from \eqref{eq:acceleration-difference-jerk}.  A sampled audit does not
license a claim of exact continuous internal-jerk verification.

\subsection{Native, shielded, and fallback identities}

% CLAIM: C12
The executed algorithm, not merely the requested endpoint or configuration
label, defines method identity.  \emph{\OrdinaryRuckig{}} executes an
unshielded native prefix; a native failure is exposed rather than silently
replaced.  \emph{\ViabilityShieldedRuckig{}} explicitly permits terminal and
next-step checks to replace a native candidate with a one-step command.
\emph{\DirectConstantJerkExecution{}} executes the governor's single segment
without secondary \Ruckig{} shaping.

% CLAIM: C12
Each cycle therefore records the declared method semantics, native follower,
whether the native command was executed, the actual command-generating
algorithm, whether a shield was requested or applied, and any fallback
controller and reason.  If a shield or fallback replaces the candidate with a
different algorithm, the resulting aggregate is not a pure native-method
comparison.  This taxonomy is an audit requirement, not a ranking of tracking
performance.
